\section{Training Objective}

\subsection{Current Baseline Objective}

The current baseline predicts both mean trajectory increments and diagonal uncertainty.
After rollout, the baseline objective is a diagonal Gaussian negative log-likelihood:
\begin{equation}
\mathcal{L}_{\mathrm{NLL}}
=
\sum_{j=1}^{H}
\sum_{d=1}^{D_{\mathrm{out}}}
\left[
\frac{\left(y_{t+j,d} - \hat{y}_{t+j,d}\right)^2}{\hat{\sigma}_{t+j,d}^{2}}
+
\log \hat{\sigma}_{t+j,d}^{2}
\right].
\end{equation}

\subsection{Dense and Sparse Supervision}

The engineering challenge is that supervision is not uniformly dense across modalities:
\begin{itemize}
  \item IMU-related target components are available on the main aligned axis,
  \item DVL-related velocity information is lower-rate and naturally sparse.
\end{itemize}

This motivates a mask-aware formulation:
\begin{equation}
\mathcal{L}_{\mathrm{mask}}
=
\sum_{j=1}^{H}
\sum_{d=1}^{D_{\mathrm{out}}}
m_{t+j,d}
\left[
\frac{\left(y_{t+j,d} - \hat{y}_{t+j,d}\right)^2}{\hat{\sigma}_{t+j,d}^{2}}
+
\log \hat{\sigma}_{t+j,d}^{2}
\right],
\end{equation}
where $m_{t+j,d} \in \{0,1\}$ encodes whether the target dimension is valid at that time step.

\subsection{Interpretation}

In the current repository, the dense baseline is already implemented, while the fully mask-aware loss is a planned extension.
The mathematical form above is therefore both:
\begin{itemize}
  \item a precise statement of the intended sparse-supervision objective,
  \item a bridge between the current baseline and future mask-aware training.
\end{itemize}
